\honest{Belief in Self-Deception}{Eliezer Yudkowsky}{2009}

I return to yesterday's woman. I put the Litany of Tarski to her: if it were known that God does not exist, should she still believe in God? She hesitated: ``So it's a counterfactual question...'' At the time I thought she could not bear to picture a world without God. Now ``I suspect'' God was not in her model of the world at all, only her belief in God. \nb{This reading rests on one hesitation over one question.}

Then she offered a counterexample to the litany: ``I believe that people are nicer than they really are.'' I explained that saying ``I believe people are nice'' means you believe you believe it, and I quoted Wittgenstein: ``If there were a verb meaning `to believe falsely', it would not have any significant first person, present indicative.'' \nb{Her sentence has the form of Moore's paradox, which the author names three days later. The Wittgenstein line comes from a section on that paradox, which says that ``I believe that this is the case'' is ``used like the assertion `This is the case'.'' That is a different view from the author's gloss.} She smiled, said it again, and added: ``I just thought I should put it that way for you.'' I do not ask what she meant by that.

I cannot imagine what it is like to be her. ``This is why intelligent people only have a certain amount of time'' to become atheists, I say. My evidence is this one conversation.

She told me that nothing should be done differently if there were definitely no God. Her religion ``seems'' to be the worship of worship; in the next sentence, she ``now believes that belief in God will save her.'' Asked whether she is always surprised when people fall short, she paused for a long time. From the pause, ``I now realize that the whole essence of her philosophy was her belief that she had deceived herself.'' I conclude: ``The attempt failed, but she is honestly unaware of this.''

This explains why I argue that self-deception is hard. You ``can't just choose to believe the sky is green,'' but you can believe you have, and then you get a placebo benefit. Yesterday that benefit came with ``I suppose''; today it does not. So by explaining how hard self-deception is, I am ``taking direct aim at the placebo benefits.''

I end by asking whether theists will now defend belief in belief, and belief in belief in belief.
